\subsection{Description of an Isotypical Module}

As in the previous subsection, A denotes a ring, S a simple left A-module, and D the field \(\operatorname{End}_{\mathrm{A}}(\mathrm{S})^{\circ}\) . We view S as an (A,D)-bimodule.

DEFINITION 3. --- Let M be an isotypical A-module of type S. A description of M with respect to S is a pair \((V, \alpha)\) , where V is a left vector space over the field D and \(\alpha : S \otimes_{D} V \to M\) is an isomorphism of A-modules.

Every isotypical A-module M of type S has a canonical description: it is the pair \((\mathrm{Hom}_{\mathrm{A}}(\mathrm{S},\mathrm{M}),\alpha_{\mathrm{M}})\) , where \(\alpha_{M}:S\otimes_{D}\mathrm{Hom}_{\mathrm{A}}(\mathrm{S},\mathrm{M})\to\mathrm{M}\) is the isomorphism of A-modules characterized by \(\alpha_{\mathrm{M}}(s\otimes f)=f(s)\) (VIII, p.~62, Proposition 3, a)).

THEOREM 2. --- Let M be an isotypical A-module of type S and \((V,\alpha)\) a description of M. Denote by \(\mathcal{D}_{\mathrm{D}}(\mathrm{V})\) the set, ordered by inclusion, of D-linear subspaces of V, and by \(\mathcal{D}_{\mathrm{A}}(\mathrm{M})\) that of A-submodules of M. For every \(W \in \mathcal{D}_{\mathrm{D}}(\mathrm{V})\) , identify the A-module \(S \otimes_{D} W\) with its canonical image in \(S \otimes_{D} V\) .

\begin{enumerate}
\def\labelenumi{\alph{enumi})}
\item
  The mapping \(W \mapsto \alpha(S \otimes_{D} W)\) is an isomorphism of ordered sets from \(\mathcal{D}_{\mathrm{D}}(\mathrm{V})\) to \(\mathcal{D}_{\mathrm{A}}(\mathrm{M})\) .
\item
  The inverse isomorphism sends a submodule N of M to the linear subspace of V consisting of the elements v such that \(\alpha(s \otimes v)\) belongs to N for every \(s \in S\) .
\end{enumerate}

For \(\mathrm{W} \in \mathcal{D}_{\mathrm{D}}(\mathrm{V})\) , set \(\varphi(\mathrm{W}) = \alpha(\mathrm{S} \otimes_{\mathrm{D}} \mathrm{W})\) . For \(\mathrm{N} \in \mathcal{D}_{\mathrm{A}}(\mathrm{M})\) , denote by \(\psi(\mathrm{N})\) the set of elements \(v \in \mathrm{V}\) such that \(\alpha(s \otimes v) \in \mathrm{N}\) for every \(s \in \mathrm{S}\) . This defines two mappings \(\varphi: \mathcal{D}_{\mathrm{D}}(\mathrm{V}) \to \mathcal{D}_{\mathrm{A}}(\mathrm{M})\) and \(\psi: \mathcal{D}_{\mathrm{A}}(\mathrm{M}) \to \mathcal{D}_{\mathrm{D}}(\mathrm{V})\) . They are clearly increasing.

Let N be a submodule of M. It is isotypical of type S (VIII, p.~61, Proposition 2). Set \(W = \psi(N)\) . By Proposition 3, b) of VIII, p.~62, the A-linear mappings \(h : S \to M\) are simply the mappings \(s \mapsto \alpha(s \otimes v)\) , where v runs through V. Those with image contained in N are the mappings \(s \mapsto \alpha(s \otimes w)\) where w runs through W; their images generate N because N is isotypical of type S. We therefore have \(\alpha(S \otimes_{D} W) = N\) , that is, \(\varphi(\psi(N)) = N\) . This proves that \(\varphi \circ \psi\) is the identity mapping on \(\mathcal{D}_{\mathrm{A}}(\mathrm{M})\) . In particular, \(\varphi\) is surjective and \(\psi\) is injective.

To finish the proof, it suffices to prove that the mapping \(\varphi\) is injective. Let W and \(W'\) be linear subspaces of V such that \(\varphi(\mathrm{W}) = \varphi(\mathrm{W}')\) . The vector spaces \(S \otimes_{D} W\) and \(S \otimes_{D} W'\) coincide when viewed as linear subspaces of \(S \otimes_{D} V\) . Choose a nonzero linear form f on the D-vector space S, and let \(g : S \otimes_{D} V \to V\) be the group homomorphism defined by \(g(s \otimes v) = f(s)v\) . We have \(W = g(S \otimes_{D} W) = g(S \otimes_{D} W') = W'\) , so \(\varphi\) is injective.

Remark 1. --- Let M be an isotypical A-module of type S and \((V, \alpha)\) a description of M. Then M has finite length if and only if V is finite-dimensional, and in this case, we have the relation

\[
\mathrm{long} _ {\mathrm{A}} (\mathrm{M}) = \mathrm{dim} _ {\mathrm{D}} (\mathrm{V}).
\]

COROLLARY 1. --- Let M be an isotypical A-module of type S. For every A-submodule N of M, identify \(\operatorname{Hom}_{\mathrm{A}}(\mathrm{S},\mathrm{N})\) with the linear subspace of \(\operatorname{Hom}_{\mathrm{A}}(\mathrm{S},\mathrm{M})\) consisting of the mappings with image contained in N.

\begin{enumerate}
\def\labelenumi{\alph{enumi})}
\item
  The mapping \(N \mapsto \operatorname{Hom}_{\mathrm{A}}(\mathrm{S}, \mathrm{N})\) is an isomorphism of ordered sets from \(\mathcal{D}_{\mathrm{A}}(\mathrm{M})\) to \(\mathcal{D}_{\mathrm{D}}(\operatorname{Hom}_{\mathrm{A}}(\mathrm{S}, \mathrm{M}))\) .
\item
  The inverse bijection sends a linear subspace \(\mathrm{W}\) of \(\operatorname{Hom}_{\mathrm{A}}(\mathrm{S}, \mathrm{M})\) to the submodule \(\sum_{h \in \mathrm{W}} h(\mathrm{S})\) of \(\mathrm{M}\) .
\end{enumerate}

This is a reformulation of Theorem 2 when we take \((\mathrm{V},\alpha)\) to be the canonical description of M.

COROLLARY 2. --- Let V be a left vector space over D and F a set of endomorphisms of V. An A-submodule of \(S \otimes_{D} V\) is stable under all endomorphisms

\(1_{\mathrm{S}} \otimes u\) , where \(u\) runs through \(\mathcal{F}\) , if and only if it is of the form \(\mathrm{S} \otimes_{\mathrm{D}} \mathrm{W}\) , where \(\mathrm{W}\) is a linear subspace of \(\mathrm{V}\) that is stable under all endomorphisms belonging to \(\mathcal{F}\) .

Indeed, by Theorem 2, every A-submodule N of \(S \otimes_{D} V\) is equal to \(S \otimes_{D} W\) , where W is the linear subspace of V consisting of the elements v such that \(s \otimes v\) belongs to N for every \(s \in S\) .

THEOREM 3. --- Let M and M' be isotypical A-modules of type S. Let \((V, \alpha)\) and \((V', \alpha')\) be descriptions of M and M', respectively. For any D-linear mapping \(f : V \to V'\) , denote by \(\widetilde{f} : M \to M'\) the unique A-linear mapping that makes the following diagram commute:

\[
\begin{array}{c} \mathrm{S} \otimes_ {\mathrm{D}} \mathrm{V} \xrightarrow {\alpha} \mathrm{M} \\ \Biggl \downarrow_ {1 _ {\mathrm{S}} \otimes f} \\ \mathrm{S} \otimes_ {\mathrm{D}} \mathrm{V} ^ {\prime} \xrightarrow {\alpha^ {\prime}} \mathrm{M} ^ {\prime}. \end{array}\tag{9}
\]

The mapping \(f \mapsto \widetilde{f}\) from \(\mathrm{Hom}_{\mathrm{D}}(\mathrm{V},\mathrm{V}')\) to \(\mathrm{Hom}_{\mathrm{A}}(\mathrm{M},\mathrm{M}')\) is a group isomorphism.

It suffices to prove that the Z-linear mapping \(u \mapsto 1_{S} \otimes u\) from \(\mathrm{Hom}_{\mathrm{D}}(\mathrm{V}, \mathrm{V}')\) to \(\mathrm{Hom}_{\mathrm{A}}(\mathrm{S} \otimes_{\mathrm{D}} \mathrm{V}, \mathrm{S} \otimes_{\mathrm{D}} \mathrm{V}')\) is bijective. In the notation of No.3 applied to the (A,D)-bimodule S, this corresponds to showing that the mapping

\[
\mathcal {T}: \mathrm{Hom} _ {\mathrm{D}} (\mathrm{V}, \mathrm{V} ^ {\prime}) \longrightarrow \mathrm{Hom} _ {\mathrm{A}} (\mathcal {T} (\mathrm{V}), \mathcal {T} (\mathrm{V} ^ {\prime}))
\]

is bijective. But by Remark 1 of VIII, p.~60, since the adjunction isomorphism (VIII, p.~59) is bijective, this corresponds to showing that the mapping that sends u to \(\beta_{V'} \circ u\) is bijective, which follows from the fact that the D-linear mapping \(\beta_{V'}\) is bijective (VIII, p.~62, Proposition 3, b)).

We keep the notation of Theorem 3. Let \(\mathrm{M}''\) be an isotypical A-module of type S, and let \((\mathrm{V}'',\alpha '')\) be a description of \(\mathrm{M}''\) . For every \(f\in \mathrm{Hom}_{\mathrm{D}}(\mathrm{V},\mathrm{V}')\) and every \(g\in \mathrm{Hom}_{\mathrm{D}}(\mathrm{V}',\mathrm{V}'')\) , we have \(\widetilde{g\circ f} = \widetilde{g}\circ \widetilde{f}\) . In particular, for \(\mathrm{M} = \mathrm{M}'\) , \(\mathrm{V} = \mathrm{V}'\) , and \(\alpha = \alpha '\) , the mapping \(f\mapsto \widetilde{f}\) from \(\mathrm{End}_{\mathrm{D}}(\mathrm{V})\) to \(\mathrm{End}_{\mathrm{A}}(\mathrm{M})\) is a ring isomorphism.

Remark 2. --- Let M be an isotypical A-module of type S, and let \((V,\alpha)\) be a description of M. Let B be a subring of the ring \(\mathrm{End}_{\mathrm{A}}(\mathrm{M})^{\circ}\) . The ring isomorphism from \(\mathrm{End}_{\mathrm{D}}(\mathrm{V})^{\circ}\) to \(\mathrm{End}_{\mathrm{A}}(\mathrm{M})^{\circ}\) endows V with the structure of a \((D,B)\) -bimodule, so that \(\alpha\) is an isomorphism of \((A,B)\) -bimodules. There exists an isomorphism from the set of \((D,B)\) -sub-bimodules of V, ordered by inclusion, to the set of (A,B)-sub-bimodules of M (VIII, p.~62, Theorem 2 and VIII, p.~63, Corollary 2).

COROLLARY. --- Let M and M' be isotypical A-modules of type S. The mapping \(u \mapsto \operatorname{Hom}(1_{\mathrm{S}}, u)\) from \(\operatorname{Hom}_{\mathrm{A}}(\mathrm{M}, \mathrm{M}')\) to \(\operatorname{Hom}_{\mathrm{D}}(\operatorname{Hom}_{\mathrm{A}}(\mathrm{S}, \mathrm{M}), \operatorname{Hom}_{\mathrm{A}}(\mathrm{S}, \mathrm{M}'))\) is a group isomorphism. When M is equal to \(M'\) , it is a ring isomorphism from \(\operatorname{End}_{\mathrm{A}}(\mathrm{M})\) to \(\operatorname{End}_{\mathrm{D}}(\operatorname{Hom}_{\mathrm{A}}(\mathrm{S}, \mathrm{M}))\) .

Because of the commutativity of diagram (I) of VIII, p.~59, this corollary follows from Theorem 3, applied to the canonical descriptions of M and M'.
% semantic-fragments: legacy-input-seam
\relax{}
